\documentclass[11pt,a4paper]{article}
\usepackage[utf8]{inputenc}
\usepackage[T1]{fontenc}
\usepackage[catalan]{babel}
\usepackage{lmodern}
\usepackage[sfdefault,lining]{FiraSans}
\usepackage[a4paper,top=2.3cm,bottom=2.2cm,left=1.8cm,right=1.8cm,headheight=1.3cm,headsep=0.4cm,footskip=1cm]{geometry}
\usepackage{amsmath,amssymb}
\usepackage{icomma}
\usepackage{graphicx}
\usepackage[table]{xcolor}
\usepackage{tikz}
\usepackage{fontawesome5}
\usepackage{ccicons}
\usepackage{fancyhdr}
\usepackage{lastpage}
\usepackage{tasks}
\usepackage{needspace}
\usepackage[most]{tcolorbox}
\usepackage{array}
\usepackage[hidelinks,pdftitle={Fitxa d'exercicis de la unitat 1: Nombres reals. 4t d'ESO Matemàtiques B},pdfsubject={Matemàtiques B 4t ESO},pdfkeywords={nombres reals, intervals, semirectes, unió, intersecció, radicals, racionalització, logaritmes}]{hyperref}

\definecolor{navy}{HTML}{1B2748}
\definecolor{accent}{HTML}{E8603A}
\definecolor{teal}{HTML}{1F9E90}
\definecolor{soft}{HTML}{EEF1F8}
\definecolor{muted}{HTML}{5D6685}

\setlength{\parindent}{0pt}
\setlength{\parskip}{3pt}
\settasks{label=\alph*), label-format=\bfseries\color{navy}, label-width=1.3em, item-indent=1.8em, column-sep=0.6em, after-item-skip=5pt, before-skip=3pt, after-skip=4pt}

% ---------- capçalera i peu ----------
\newcommand{\llicencia}{\href{https://creativecommons.org/licenses/by-nc-sa/4.0/deed.ca}{\ccbyncsa}\ \ Llicència Creative Commons Reconeixement-NoComercial-CompartirIgual 4.0 Internacional (CC BY-NC-SA 4.0)}
\pagestyle{fancy}
\fancyhf{}
\fancyhead[L]{\raisebox{-0.35\height}{\includegraphics[height=0.95cm]{logo-gva.png}}}
\fancyhead[C]{\small\color{muted}Fitxa d'exercicis · Unitat 1: Nombres reals · 4t ESO B}
\fancyhead[R]{\raisebox{-0.35\height}{\includegraphics[height=0.95cm]{logo-vermellar.png}}}
\fancyfoot[L]{\scriptsize\color{muted}\llicencia}
\fancyfoot[R]{\small\color{muted}\thepage}
\renewcommand{\headrulewidth}{0.4pt}
\renewcommand{\headrule}{\color{muted!40}\hrule height \headrulewidth}
\fancypagestyle{primera}{\fancyhf{}\renewcommand{\headrulewidth}{0pt}
  \fancyfoot[L]{\scriptsize\color{muted}\llicencia}
  \fancyfoot[R]{\small\color{muted}\thepage}}

% ---------- nivells i exercicis ----------
\newcommand{\estrelles}[1]{{\color{accent}\ifcase#1\or\faStar\faStar[regular]\faStar[regular]\or\faStar\faStar\faStar[regular]\or\faStar\faStar\faStar\fi}}
\newcounter{exer}
\newcommand{\ex}[1]{\par\needspace{4\baselineskip}\medskip\refstepcounter{exer}\noindent{\color{navy}\bfseries\theexer.}\ \ {\small\estrelles{#1}}\ \ }

\newcommand{\seccio}[2]{\par\needspace{16\baselineskip}\bigskip
  \begin{tcolorbox}[colback=navy,colframe=navy,coltext=white,arc=2mm,boxrule=0pt,left=3mm,right=3mm,top=1.2mm,bottom=1.2mm]
  {\large\bfseries #1\quad #2}
  \end{tcolorbox}\vspace{-2pt}}

\newcommand{\exf}[1]{\par\medskip\refstepcounter{exer}\noindent{\color{navy}\bfseries\theexer.}\ \ {\small\estrelles{#1}}\ \ }
\newcolumntype{C}[1]{>{\centering\arraybackslash}p{#1}}
% quadre de recordatori (només a la fitxa d'exercicis)
\newcommand{\recorda}[1]{\par\smallskip\begin{tcolorbox}[colback=accent!7,colframe=accent!55,arc=2mm,boxrule=0.5pt,left=2.5mm,right=2.5mm,top=1mm,bottom=1mm]\small{\color{accent}\faLightbulb[regular]}\ \textbf{Recorda:}\ #1\end{tcolorbox}}
\newcommand{\res}[1]{#1}\newcommand{\rest}[1]{#1}

\def\fmneed{11}% que el títol d'una secció no quede sol al final del full
% =====================================================================
% Capçalera FentMates de les fitxes: un color per a cada curs.
% Els colors són els de la pàgina inicial (LEVELS a assets/js/materials.js);
% el logo ve de _eines/logo/pdf/ (python3 _eines/logo/genera.py).
%
% Ús, al final del preàmbul de la unitat:
%   % =====================================================================
% Capçalera FentMates de les fitxes: un color per a cada curs.
% Els colors són els de la pàgina inicial (LEVELS a assets/js/materials.js);
% el logo ve de _eines/logo/pdf/ (python3 _eines/logo/genera.py).
%
% Ús, al final del preàmbul de la unitat:
%   % =====================================================================
% Capçalera FentMates de les fitxes: un color per a cada curs.
% Els colors són els de la pàgina inicial (LEVELS a assets/js/materials.js);
% el logo ve de _eines/logo/pdf/ (python3 _eines/logo/genera.py).
%
% Ús, al final del preàmbul de la unitat:
%   \input{../comu/fentmates.tex}
%   \fmcurs{1eso}
% i, al començament del document:
%   \capcalera{Matemàtiques}{1r ESO}{Fitxa d'exercicis · Unitat 2}{Divisibilitat}
% Opcional abans de \input: \def\fmneed{14} (espai mínim abans de cada secció).
% =====================================================================
\makeatletter
\def\fm@logodir{../../logo/pdf}
\providecommand{\fmneed}{10}
\def\fm@cursnom{}

% clau del curs, color i color del text damunt del color
\newcommand{\fm@curs}[3]{\@namedef{fm@color@#1}{#2}\@namedef{fm@text@#1}{#3}}
\fm@curs{1eso}{E6B84A}{navy}
\fm@curs{2eso}{3B82F6}{white}
\fm@curs{3eso}{FF7A50}{white}
\fm@curs{4eso-a}{A78BFA}{white}
\fm@curs{4eso-b}{6D28D9}{white}
\fm@curs{1batx-mi}{9B1C1C}{white}
\fm@curs{1batx-ccss}{EF6B6B}{white}
\fm@curs{2batx-mii}{1E6B47}{white}
\fm@curs{2batx-ccss}{2BB3A3}{white}

\newcommand{\fmcurs}[1]{%
  \@ifundefined{fm@color@#1}{\PackageError{fentmates}{Curs desconegut: #1}{}}{}%
  \edef\fm@tmp{\noexpand\definecolor{curs}{HTML}{\@nameuse{fm@color@#1}}%
               \noexpand\colorlet{curstext}{\@nameuse{fm@text@#1}}}\fm@tmp
  \colorlet{cursfosc}{curs!72!black}%
  \def\fm@logo{\fm@logodir/vertical-#1.pdf}%
  \def\fm@marca{\fm@logodir/marca-#1.pdf}}

% ---------- capçalera de la primera pàgina ----------
% #1 matèria (petit), #2 curs, #3 tipus de document i unitat, #4 títol
\newcommand{\capcalera}[4]{%
  \gdef\fm@cursnom{#2}%
  \par\noindent
  \begin{tcolorbox}[colback=curs,colframe=curs,boxrule=0pt,arc=2.2mm,width=3.5cm,height=2.15cm,
      valign=center,left=3mm,right=2mm,top=1mm,bottom=1mm,nobeforeafter,box align=center]
    \color{curstext}{\footnotesize #1}\par\vspace{1pt}{\fontsize{19}{21}\selectfont\bfseries #2\par}%
  \end{tcolorbox}\hspace{5mm}%
  \begin{minipage}[c]{\dimexpr\linewidth-3.5cm-5mm-2.9cm\relax}
    {\color{cursfosc}\small\bfseries\MakeUppercase{#3}}\par\vspace{2pt}
    {\color{navy}\fontsize{20}{23}\selectfont\bfseries #4\par}
  \end{minipage}\hfill
  \begin{minipage}[c]{2.6cm}\raggedleft\includegraphics[width=2.6cm]{\fm@logo}\end{minipage}\par
  \vspace{2pt}}

% ---------- barres de secció del color del curs ----------
\renewcommand{\seccio}[2]{\par\needspace{\fmneed\baselineskip}\bigskip
  \begin{tcolorbox}[colback=curs,colframe=curs,coltext=curstext,arc=2mm,boxrule=0pt,left=3mm,right=3mm,top=1.2mm,bottom=1.2mm]
  {\large\bfseries #1\quad #2}
  \end{tcolorbox}\vspace{-2pt}}

% ---------- estrelles de dificultat del color del curs (a Batxillerat no se'n posen) ----------
\@ifundefined{estrelles}{}{%
  \renewcommand{\estrelles}[1]{{\color{curs}\ifcase#1\or\faStar\faStar[regular]\faStar[regular]\or\faStar\faStar\faStar[regular]\or\faStar\faStar\faStar\fi}}}

% ---------- capçalera i peu de les pàgines ----------
% primera pàgina: logos institucionals al peu
\fancypagestyle{primera}{\fancyhf{}\renewcommand{\headrulewidth}{0pt}
  \fancyfoot[L]{\raisebox{-0.32\height}{\includegraphics[height=0.78cm]{logo-gva.png}}\hspace{5mm}%
    \raisebox{-0.32\height}{\includegraphics[height=0.78cm]{logo-vermellar.png}}}
  \fancyfoot[R]{\scriptsize\color{muted}\href{https://creativecommons.org/licenses/by-nc-sa/4.0/deed.ca}{\ccbyncsa}\ \ CC BY-NC-SA 4.0 · fentmates.org\qquad\small\thepage}}
% la resta: el símbol a l'esquerra i el curs a la dreta, amb el títol al centre com abans
\fancyhead[L]{\raisebox{-0.3\height}{\includegraphics[height=0.62cm]{\fm@marca}}}
\fancyhead[R]{\small\bfseries\color{cursfosc}\fm@cursnom}
\renewcommand{\headrule}{\color{curs!45}\hrule height \headrulewidth}
\makeatother

%   \fmcurs{1eso}
% i, al començament del document:
%   \capcalera{Matemàtiques}{1r ESO}{Fitxa d'exercicis · Unitat 2}{Divisibilitat}
% Opcional abans de \input: \def\fmneed{14} (espai mínim abans de cada secció).
% =====================================================================
\makeatletter
\def\fm@logodir{../../logo/pdf}
\providecommand{\fmneed}{10}
\def\fm@cursnom{}

% clau del curs, color i color del text damunt del color
\newcommand{\fm@curs}[3]{\@namedef{fm@color@#1}{#2}\@namedef{fm@text@#1}{#3}}
\fm@curs{1eso}{E6B84A}{navy}
\fm@curs{2eso}{3B82F6}{white}
\fm@curs{3eso}{FF7A50}{white}
\fm@curs{4eso-a}{A78BFA}{white}
\fm@curs{4eso-b}{6D28D9}{white}
\fm@curs{1batx-mi}{9B1C1C}{white}
\fm@curs{1batx-ccss}{EF6B6B}{white}
\fm@curs{2batx-mii}{1E6B47}{white}
\fm@curs{2batx-ccss}{2BB3A3}{white}

\newcommand{\fmcurs}[1]{%
  \@ifundefined{fm@color@#1}{\PackageError{fentmates}{Curs desconegut: #1}{}}{}%
  \edef\fm@tmp{\noexpand\definecolor{curs}{HTML}{\@nameuse{fm@color@#1}}%
               \noexpand\colorlet{curstext}{\@nameuse{fm@text@#1}}}\fm@tmp
  \colorlet{cursfosc}{curs!72!black}%
  \def\fm@logo{\fm@logodir/vertical-#1.pdf}%
  \def\fm@marca{\fm@logodir/marca-#1.pdf}}

% ---------- capçalera de la primera pàgina ----------
% #1 matèria (petit), #2 curs, #3 tipus de document i unitat, #4 títol
\newcommand{\capcalera}[4]{%
  \gdef\fm@cursnom{#2}%
  \par\noindent
  \begin{tcolorbox}[colback=curs,colframe=curs,boxrule=0pt,arc=2.2mm,width=3.5cm,height=2.15cm,
      valign=center,left=3mm,right=2mm,top=1mm,bottom=1mm,nobeforeafter,box align=center]
    \color{curstext}{\footnotesize #1}\par\vspace{1pt}{\fontsize{19}{21}\selectfont\bfseries #2\par}%
  \end{tcolorbox}\hspace{5mm}%
  \begin{minipage}[c]{\dimexpr\linewidth-3.5cm-5mm-2.9cm\relax}
    {\color{cursfosc}\small\bfseries\MakeUppercase{#3}}\par\vspace{2pt}
    {\color{navy}\fontsize{20}{23}\selectfont\bfseries #4\par}
  \end{minipage}\hfill
  \begin{minipage}[c]{2.6cm}\raggedleft\includegraphics[width=2.6cm]{\fm@logo}\end{minipage}\par
  \vspace{2pt}}

% ---------- barres de secció del color del curs ----------
\renewcommand{\seccio}[2]{\par\needspace{\fmneed\baselineskip}\bigskip
  \begin{tcolorbox}[colback=curs,colframe=curs,coltext=curstext,arc=2mm,boxrule=0pt,left=3mm,right=3mm,top=1.2mm,bottom=1.2mm]
  {\large\bfseries #1\quad #2}
  \end{tcolorbox}\vspace{-2pt}}

% ---------- estrelles de dificultat del color del curs (a Batxillerat no se'n posen) ----------
\@ifundefined{estrelles}{}{%
  \renewcommand{\estrelles}[1]{{\color{curs}\ifcase#1\or\faStar\faStar[regular]\faStar[regular]\or\faStar\faStar\faStar[regular]\or\faStar\faStar\faStar\fi}}}

% ---------- capçalera i peu de les pàgines ----------
% primera pàgina: logos institucionals al peu
\fancypagestyle{primera}{\fancyhf{}\renewcommand{\headrulewidth}{0pt}
  \fancyfoot[L]{\raisebox{-0.32\height}{\includegraphics[height=0.78cm]{logo-gva.png}}\hspace{5mm}%
    \raisebox{-0.32\height}{\includegraphics[height=0.78cm]{logo-vermellar.png}}}
  \fancyfoot[R]{\scriptsize\color{muted}\href{https://creativecommons.org/licenses/by-nc-sa/4.0/deed.ca}{\ccbyncsa}\ \ CC BY-NC-SA 4.0 · fentmates.org\qquad\small\thepage}}
% la resta: el símbol a l'esquerra i el curs a la dreta, amb el títol al centre com abans
\fancyhead[L]{\raisebox{-0.3\height}{\includegraphics[height=0.62cm]{\fm@marca}}}
\fancyhead[R]{\small\bfseries\color{cursfosc}\fm@cursnom}
\renewcommand{\headrule}{\color{curs!45}\hrule height \headrulewidth}
\makeatother

%   \fmcurs{1eso}
% i, al començament del document:
%   \capcalera{Matemàtiques}{1r ESO}{Fitxa d'exercicis · Unitat 2}{Divisibilitat}
% Opcional abans de \input: \def\fmneed{14} (espai mínim abans de cada secció).
% =====================================================================
\makeatletter
\def\fm@logodir{../../logo/pdf}
\providecommand{\fmneed}{10}
\def\fm@cursnom{}

% clau del curs, color i color del text damunt del color
\newcommand{\fm@curs}[3]{\@namedef{fm@color@#1}{#2}\@namedef{fm@text@#1}{#3}}
\fm@curs{1eso}{E6B84A}{navy}
\fm@curs{2eso}{3B82F6}{white}
\fm@curs{3eso}{FF7A50}{white}
\fm@curs{4eso-a}{A78BFA}{white}
\fm@curs{4eso-b}{6D28D9}{white}
\fm@curs{1batx-mi}{9B1C1C}{white}
\fm@curs{1batx-ccss}{EF6B6B}{white}
\fm@curs{2batx-mii}{1E6B47}{white}
\fm@curs{2batx-ccss}{2BB3A3}{white}

\newcommand{\fmcurs}[1]{%
  \@ifundefined{fm@color@#1}{\PackageError{fentmates}{Curs desconegut: #1}{}}{}%
  \edef\fm@tmp{\noexpand\definecolor{curs}{HTML}{\@nameuse{fm@color@#1}}%
               \noexpand\colorlet{curstext}{\@nameuse{fm@text@#1}}}\fm@tmp
  \colorlet{cursfosc}{curs!72!black}%
  \def\fm@logo{\fm@logodir/vertical-#1.pdf}%
  \def\fm@marca{\fm@logodir/marca-#1.pdf}}

% ---------- capçalera de la primera pàgina ----------
% #1 matèria (petit), #2 curs, #3 tipus de document i unitat, #4 títol
\newcommand{\capcalera}[4]{%
  \gdef\fm@cursnom{#2}%
  \par\noindent
  \begin{tcolorbox}[colback=curs,colframe=curs,boxrule=0pt,arc=2.2mm,width=3.5cm,height=2.15cm,
      valign=center,left=3mm,right=2mm,top=1mm,bottom=1mm,nobeforeafter,box align=center]
    \color{curstext}{\footnotesize #1}\par\vspace{1pt}{\fontsize{19}{21}\selectfont\bfseries #2\par}%
  \end{tcolorbox}\hspace{5mm}%
  \begin{minipage}[c]{\dimexpr\linewidth-3.5cm-5mm-2.9cm\relax}
    {\color{cursfosc}\small\bfseries\MakeUppercase{#3}}\par\vspace{2pt}
    {\color{navy}\fontsize{20}{23}\selectfont\bfseries #4\par}
  \end{minipage}\hfill
  \begin{minipage}[c]{2.6cm}\raggedleft\includegraphics[width=2.6cm]{\fm@logo}\end{minipage}\par
  \vspace{2pt}}

% ---------- barres de secció del color del curs ----------
\renewcommand{\seccio}[2]{\par\needspace{\fmneed\baselineskip}\bigskip
  \begin{tcolorbox}[colback=curs,colframe=curs,coltext=curstext,arc=2mm,boxrule=0pt,left=3mm,right=3mm,top=1.2mm,bottom=1.2mm]
  {\large\bfseries #1\quad #2}
  \end{tcolorbox}\vspace{-2pt}}

% ---------- estrelles de dificultat del color del curs (a Batxillerat no se'n posen) ----------
\@ifundefined{estrelles}{}{%
  \renewcommand{\estrelles}[1]{{\color{curs}\ifcase#1\or\faStar\faStar[regular]\faStar[regular]\or\faStar\faStar\faStar[regular]\or\faStar\faStar\faStar\fi}}}

% ---------- capçalera i peu de les pàgines ----------
% primera pàgina: logos institucionals al peu
\fancypagestyle{primera}{\fancyhf{}\renewcommand{\headrulewidth}{0pt}
  \fancyfoot[L]{\raisebox{-0.32\height}{\includegraphics[height=0.78cm]{logo-gva.png}}\hspace{5mm}%
    \raisebox{-0.32\height}{\includegraphics[height=0.78cm]{logo-vermellar.png}}}
  \fancyfoot[R]{\scriptsize\color{muted}\href{https://creativecommons.org/licenses/by-nc-sa/4.0/deed.ca}{\ccbyncsa}\ \ CC BY-NC-SA 4.0 · fentmates.org\qquad\small\thepage}}
% la resta: el símbol a l'esquerra i el curs a la dreta, amb el títol al centre com abans
\fancyhead[L]{\raisebox{-0.3\height}{\includegraphics[height=0.62cm]{\fm@marca}}}
\fancyhead[R]{\small\bfseries\color{cursfosc}\fm@cursnom}
\renewcommand{\headrule}{\color{curs!45}\hrule height \headrulewidth}
\makeatother

\fmcurs{4eso-b}

\begin{document}
\thispagestyle{primera}

\capcalera{Matemàtiques B}{4t ESO}{Fitxa d'exercicis · Unitat 1}{\raggedright Nombres reals}
\vspace{10pt}
Nom i cognoms: \rule{8.4cm}{0.4pt}\hfill Grup: \rule{1.6cm}{0.4pt}\hfill Data: \rule{2.2cm}{0.4pt}
\vspace{8pt}

\begin{tcolorbox}[colback=soft,colframe=soft,arc=2mm,left=3mm,right=3mm,top=1.5mm,bottom=1.5mm]
\small Exercicis graduats per nivells de dificultat:\quad \estrelles{1}\ bàsic \qquad \estrelles{2}\ mitjà \qquad \estrelles{3}\ avançat.\\[2pt] Resol-los al quadern mostrant tots els passos. Dona els resultats de manera exacta (fraccions i radicals simplificats, sense arrels al denominador). Fes servir la calculadora només quan l'enunciat ho diga.
\end{tcolorbox}

% =====================================================
\seccio{1}{Classificació dels nombres}

\exf{1} Indica el menor conjunt numèric ($\mathbb N$, $\mathbb Z$, $\mathbb Q$ o $\mathbb I$) al qual pertany cada nombre:\begin{tasks}(5)
  \task $\sqrt{36}$
  \task $-7$
  \task $\frac35$
  \task $\sqrt7$
  \task $0,25$
  \task $-\sqrt9$
  \task $2,\widehat3$
  \task $\pi$
  \task $\frac{12}{4}$
  \task $-\frac{18}{6}$
\end{tasks}

\ex{1} Indica si els nombres següents són racionals o irracionals. Justifica-ho:\begin{tasks}(3)
  \task $1,212121\ldots$
  \task $1,2122122212222\ldots$
  \task $\sqrt{\frac{16}{9}}$
  \task $\sqrt{20}$
  \task $\frac{\pi}{4}$
  \task $\sqrt{0,25}$
\end{tasks}

\ex{2} Troba la fracció generatriu i simplifica-la:\begin{tasks}(5)
  \task $0,\widehat7$
  \task $2,\widehat{45}$
  \task $1,3\widehat6$
  \task $0,12\widehat5$
  \task $3,\widehat9$
\end{tasks}\recorda{Al numerador, el nombre sense coma menys la part no periòdica; al denominador, tants 9 com xifres té el període i tants 0 com xifres té l'anteperíode.}

\ex{2} Raona si les afirmacions següents són verdaderes o falses:\begin{tasks}(1)
  \task Tot nombre enter és racional.
  \task Tot nombre decimal es pot escriure com una fracció.
  \task L'arrel quadrada d'un nombre natural sempre és irracional.
  \task La suma d'un racional i un irracional és irracional.
  \task $\sqrt{-16}$ és un nombre enter.
\end{tasks}


% =====================================================
\seccio{2}{Intervals i semirectes}

\exf{1} Completa la taula:

\begin{center}\renewcommand{\arraystretch}{1.5}\begin{tabular}{|C{4.6cm}|C{3.2cm}|C{6cm}|}\hline\rowcolor{soft}\textbf{Desigualtat} & \textbf{Interval} & \textbf{Representació}\\\hline$-1\le x<4$ & & \\\hline & $(-\infty,\, 2]$ & \\\hline $x>-3$ & & \\\hline & $(0,\, 5)$ & \\\hline $x<1$ & & \\\hline\end{tabular}\end{center}

\ex{1} Escriu com a interval o semirecta i representa:\begin{tasks}(3)
  \task $x>2$
  \task $-1\le x\le4$
  \task $x\le-3$
  \task $0<x<5$
  \task $-2\le x<3$
  \task $x\ge\frac12$
\end{tasks}

\ex{1} Escriu la desigualtat que compleixen els nombres de cada interval i representa'l:\begin{tasks}(4)
  \task $[-3,\, 1)$
  \task $(-\infty,\, 2]$
  \task $(4,\, +\infty)$
  \task $(-5,\, 0)$
\end{tasks}

\ex{2} Escriu com a interval o semirecta:\begin{tasks}(1)
  \task Els nombres més grans que $-2$ i, com a màxim, iguals a 6.
  \task Els nombres reals positius.
  \task Els nombres reals no negatius.
  \task Els nombres que estan entre 1 i 5, amb el 5 inclòs.
  \task Els nombres la distància dels quals al 0 és menor que 3.
\end{tasks}


% =====================================================
\seccio{3}{Unió i intersecció d'intervals}

\exf{2} Donats els intervals $A=[-3,\, 2]$, $B=(0,\, 5]$, $C=[4,\, +\infty)$ i $D=(-\infty,\, 1)$, calcula i representa:\begin{tasks}(4)
  \task $A\cup B$
  \task $A\cap B$
  \task $B\cap C$
  \task $A\cap C$
  \task $A\cup D$
  \task $B\cap D$
  \task $C\cup D$
  \task $A\cap D$
\end{tasks}\recorda{La unió $\cup$ són els nombres que estan en algun dels dos; la intersecció $\cap$, els que estan en tots dos alhora.}

\ex{2} Calcula $E\cup G$ i $E\cap G$, amb $E=(-\infty,\, 3]$ i $G=(-1,\, 6)$. Després fes el mateix amb $H=[2,\, +\infty)$ i $K=(-\infty,\, 2)$. Què observes?

\ex{1} Escriu en forma d'interval els nombres que compleixen alhora les dues desigualtats:\begin{tasks}(4)
  \task $x<3$ i $x\ge5$
  \task $x>0$ i $x<4$
  \task $x\le-1$ i $x>-6$
  \task $x<3$ i $x\ge-2$
\end{tasks}


% =====================================================
\seccio{4}{Radicals}

\exf{1} Escriu en forma de potència:\begin{tasks}(4)
  \task $\sqrt5$
  \task $\sqrt[3]{7^2}$
  \task $\sqrt[5]{x^3}$
  \task $\dfrac{1}{\sqrt3}$
\end{tasks}I en forma de radical:\begin{tasks}[start=5](4)
  \task $3^{\frac12}$
  \task $5^{\frac23}$
  \task $a^{\frac34}$
  \task $2^{-\frac12}$
\end{tasks}\recorda{$\sqrt[n]{a^m}=a^{\frac mn}$: l'exponent va al numerador i l'índex al denominador.}

\ex{1} Calcula sense calculadora:\begin{tasks}(5)
  \task $\sqrt[3]{-27}$
  \task $\sqrt[4]{16}$
  \task $\sqrt[5]{32}$
  \task $\sqrt[3]{\frac{8}{27}}$
  \task $\sqrt{0,09}$
  \task $8^{\frac23}$
  \task $16^{\frac34}$
  \task $27^{-\frac13}$
  \task $25^{\frac12}$
  \task $\sqrt{-25}$
\end{tasks}

\ex{1} Simplifica els radicals següents:\begin{tasks}(3)
  \task $\sqrt[6]{8}$
  \task $\sqrt[9]{27}$
  \task $\sqrt[4]{25}$
  \task $\sqrt[12]{a^8}$
  \task $\sqrt[10]{x^4y^6}$
  \task $\sqrt[8]{16}$
\end{tasks}

\ex{2} Redueix a índex comú i ordena de menor a major:\begin{tasks}(2)
  \task $\sqrt3$ i $\sqrt[3]5$
  \task $\sqrt2$, $\sqrt[4]5$ i $\sqrt[3]3$
\end{tasks}

\ex{2} Extrau tots els factors possibles:\begin{tasks}(3)
  \task $\sqrt{12}$
  \task $\sqrt{50}$
  \task $\sqrt{98}$
  \task $\sqrt[3]{24}$
  \task $\sqrt[3]{250}$
  \task $\sqrt[4]{32}$
  \task $\sqrt{18a^3}$
  \task $\sqrt[3]{16x^5y^3}$
  \task $\sqrt{\frac{75}{4}}$
\end{tasks}

\ex{1} Introdueix dins del radical el factor de fora:\begin{tasks}(3)
  \task $2\sqrt3$
  \task $3\sqrt2$
  \task $2\sqrt[3]5$
  \task $a\sqrt a$
  \task $2x\sqrt[3]x$
  \task $\frac12\sqrt8$
\end{tasks}

\ex{2} Fes les sumes i restes:\begin{tasks}(2)
  \task $\sqrt{12}+\sqrt{27}$
  \task $\sqrt{50}-\sqrt{18}+\sqrt8$
  \task $2\sqrt{45}-3\sqrt{20}+\sqrt5$
  \task $\sqrt[3]{16}+\sqrt[3]{54}-\sqrt[3]2$
  \task $3\sqrt8-2\sqrt{32}+\sqrt{72}$
  \task $\sqrt{75}-2\sqrt{12}-\sqrt{27}$
  \task $5\sqrt a-\sqrt{4a}+\sqrt{9a}$
  \task $\sqrt[3]{81}-2\sqrt[3]{24}+\sqrt[3]{375}$
\end{tasks}\recorda{Només es poden sumar radicals semblants (mateix índex i mateix radicand). Extrau factors abans!}

\ex{2} Calcula i simplifica:\begin{tasks}(4)
  \task $\sqrt2\cdot\sqrt8$
  \task $\sqrt3\cdot\sqrt6$
  \task $\sqrt[3]4\cdot\sqrt[3]2$
  \task $\sqrt{10}:\sqrt2$
  \task $\sqrt[3]{54}:\sqrt[3]2$
  \task $2\sqrt3\cdot5\sqrt6$
  \task $\sqrt2\cdot\sqrt[3]2$
  \task $\sqrt3:\sqrt[4]3$
\end{tasks}

\ex{2} Calcula les potències i expressa com un sol radical:\begin{tasks}(3)
  \task $(\sqrt3)^4$
  \task $(\sqrt[3]2)^6$
  \task $(2\sqrt5)^2$
  \task $(\sqrt2)^5$
  \task $(\sqrt[3]4)^2$
  \task $(\sqrt[4]8)^2$
\end{tasks}\begin{tasks}[start=7](4)
  \task $\sqrt{\sqrt5}$
  \task $\sqrt[3]{\sqrt2}$
  \task $\sqrt{2\sqrt2}$
  \task $\sqrt{3\sqrt[3]3}$
\end{tasks}

\ex{3} Opera i simplifica:\begin{tasks}(2)
  \task $(\sqrt3+\sqrt2)^2$
  \task $(\sqrt5+1)(\sqrt5-1)$
  \task $(2\sqrt3-1)^2$
  \task $\sqrt2\left(\sqrt8-\sqrt{18}\right)$
  \task $(3+\sqrt2)(1-\sqrt2)$
\end{tasks}


% =====================================================
\seccio{5}{Racionalització}

\exf{1} Racionalitza:\begin{tasks}(3)
  \task $\dfrac{1}{\sqrt2}$
  \task $\dfrac{6}{\sqrt3}$
  \task $\dfrac{5}{2\sqrt5}$
  \task $\dfrac{\sqrt2}{\sqrt6}$
  \task $\dfrac{3}{4\sqrt6}$
  \task $\dfrac{10}{\sqrt{50}}$
\end{tasks}

\ex{2} Racionalitza:\begin{tasks}(4)
  \task $\dfrac{2}{\sqrt[3]2}$
  \task $\dfrac{6}{\sqrt[3]9}$
  \task $\dfrac{1}{\sqrt[4]8}$
  \task $\dfrac{10}{\sqrt[5]4}$
\end{tasks}\recorda{Per a $\sqrt[n]{b^m}$, multipliquem dalt i baix per $\sqrt[n]{b^{\,n-m}}$, perquè l'exponent arribe a l'índex.}

\ex{2} Racionalitza multiplicant pel conjugat:\begin{tasks}(3)
  \task $\dfrac{1}{\sqrt3-1}$
  \task $\dfrac{4}{3+\sqrt5}$
  \task $\dfrac{\sqrt2}{\sqrt2+1}$
  \task $\dfrac{1+\sqrt2}{1-\sqrt2}$
  \task $\dfrac{6}{\sqrt5-\sqrt2}$
  \task $\dfrac{\sqrt3+\sqrt2}{\sqrt3-\sqrt2}$
\end{tasks}\recorda{$(a+b)(a-b)=a^2-b^2$.}

\ex{3} Racionalitza i opera:\begin{tasks}(2)
  \task $\dfrac{1}{\sqrt2-1}-\dfrac{1}{\sqrt2+1}$
  \task $\dfrac{2}{\sqrt2}+\dfrac{3}{\sqrt2+1}$
\end{tasks}


% =====================================================
\seccio{6}{Logaritmes}

\exf{1} Aplicant la definició, calcula:\begin{tasks}(5)
  \task $\log_2 32$
  \task $\log_3 81$
  \task $\log 10\,000$
  \task $\log 0,001$
  \task $\log_5\frac{1}{25}$
  \task $\log_2\frac18$
  \task $\ln e^4$
  \task $\ln\frac1e$
  \task $\log_7 7$
  \task $\log_9 1$
\end{tasks}\recorda{$\log_a b=c\iff a^c=b$. A més, $\log_a a=1$ i $\log_a 1=0$.}

\ex{2} Aplicant la definició, calcula:\begin{tasks}(4)
  \task $\log_2\sqrt2$
  \task $\log_3\sqrt[3]3$
  \task $\log_4 8$
  \task $\log_{\frac12}4$
  \task $\log_8 2$
  \task $\log_5\sqrt{125}$
  \task $\log 0,1$
  \task $\log_{\sqrt3}9$
\end{tasks}

\ex{2} Troba el valor de $x$:\begin{tasks}(3)
  \task $\log_x 49=2$
  \task $\log_2 x=5$
  \task $\log_3 x=-2$
  \task $\log_x\frac18=-3$
  \task $\log_5 x=0$
  \task $\log x=-2$
\end{tasks}

\ex{1} Calcula aplicant les propietats dels logaritmes:\begin{tasks}(3)
  \task $\log_6 2+\log_6 3$
  \task $\log_2 24-\log_2 3$
  \task $\log 25+\log 4$
  \task $\log_3 18-\log_3 2$
  \task $\log_2 16^3$
  \task $\log_5 125+\log_5 5-\log_5 25$
\end{tasks}

\ex{2} Sabent que $\log2=0,301$ i $\log3=0,477$, calcula sense calculadora:\begin{tasks}(4)
  \task $\log6$
  \task $\log8$
  \task $\log5$
  \task $\log18$
  \task $\log1,5$
  \task $\log\sqrt{12}$
  \task $\log0,2$
  \task $\log27$
\end{tasks}

\ex{2} Desenvolupa aplicant les propietats:\begin{tasks}(2)
  \task $\log(x^2y)$
  \task $\log\dfrac{100x}{y^3}$
  \task $\log\sqrt{\dfrac xy}$
  \task $\ln(e^2x)$
  \task $\log_2\dfrac{8a^3}{b}$
\end{tasks}

\ex{2} Escriu $E$ sense logaritmes:\begin{tasks}(1)
  \task $\log E=\log3+2\log x-\log y$
  \task $\log E=2-\log x$
  \task $\log E=\frac12\log a+3\log b$
\end{tasks}

\ex{1} Amb la calculadora, calcula amb tres decimals:\begin{tasks}(4)
  \task $\log50$
  \task $\ln7$
  \task $\log0,35$
  \task $\ln0,5$
  \task $\log_3 15$
  \task $\log_2 100$
  \task $\log_7 2$
  \task $\log_5 0,2$
\end{tasks}\recorda{Canvi de base: $\log_a x=\dfrac{\log x}{\log a}=\dfrac{\ln x}{\ln a}$.}

\ex{2} Calcula sense calculadora:\begin{tasks}(2)
  \task $\log1000+\ln e^3-\log_2 16$
  \task $\ln1+\log0,01$
  \task $\ln\sqrt e$
  \task $\log\sqrt[3]{100}$
\end{tasks}


% =====================================================
\seccio{7}{Problemes}

\exf{2} Un quadrat té una àrea de $50\ \text{cm}^2$. Calcula'n el costat i la diagonal en forma de radical simplificat.

\ex{2} Un rectangle fa $\sqrt{12}$ cm d'amplària i $\sqrt{27}$ cm de llargària. Calcula'n el perímetre, l'àrea i la diagonal de manera exacta.

\ex{2} Per a anar a una excursió cal tindre almenys 12 anys i menys de 17. Escriu l'interval d'edats. Si el preu reduït és per a menors de 14 anys, quines edats poden anar a l'excursió amb preu reduït?

\ex{2} Una colònia de bacteris es duplica cada hora. Quantes hores tarda a passar de 100 bacteris a un milió?

\ex{2} Ingressem 1500 € en un compte al 4\,\% d'interés compost anual. Quants anys han de passar, com a mínim, perquè hi haja 2000 €?

\ex{3} El pH d'una dissolució és $\text{pH}=-\log[\text{H}^+]$. Calcula el pH si $[\text{H}^+]=10^{-3}$ i si $[\text{H}^+]=2,5\cdot10^{-5}$. Quina de les dues és més àcida (pH menor)?

\vfill
\begin{center}\color{muted}\small\faCheckCircle[regular]\ \ Ànim! Comprova els teus resultats amb les eines interactives dels apunts de la unitat.\end{center}

\end{document}
